\makeabstract
{
Designing a high throughput screen for activators and inhibitors of the SH2-domain containing protein tyrosine phosphatases 1 (SHP1)
}
{
Morolake Orimoloye (1), Anthony C. Bishop, PhD (1)
}
{
\textsuperscript{1}Department of Chemistry, Amherst College, Amherst, MA, USA
}
{
SHP1 plays a role as both a tumor suppressor and promoter, so exploring its activity can clarify SHP1’s role in cancer development. Small molecules can be used to potentially shift the conformational equilibrium of SHP1 between the active and inactive states. However, no well-characterized allosteric activators or inhibitors currently exist. This project focuses on a mutant (S496P) that is positioned halfway between the closed and open states, exhibiting increased activity, as a tool for screening small-molecule candidates.
}
{
Wild-type SHP1 and truncated SHP1 proteins were expressed and purified. and purity was confirmed by SDS-PAGE. Enzyme kinetics were measured using a Michaelis–Menten assay to determine catalytic parameters. Inhibition assays were performed by testing five small-molecule compounds across varying concentrations to assess their effects on SHP1 activity.
}
{
Kinetic analysis revealed that full-length WT SHP1 displayed more than a threefold lower catalytic efficiency compared to truncated SHP1, highlighting the importance of the C-terminal tail. Inhibition assays produced ambiguous results: while some compounds showed reduced activity at higher concentrations, others unexpectedly increased SHP1 activity.
}
{
Although we were able to successfully express WT SHP1, the results from the inhibition assay were ambiguous, with no consistent trend observed across the compounds. Therefore, the assay will be repeated to confirm whether these molecules truly inhibit SHP1 activity. We also plan on expanding our screening to a broader library of activators and inhibitors. A particular focus will be placed on evaluating the effects on the S496P SHP1 mutant.
}

\newpage
